\chapter{麦克斯韦方程组}

电磁学在本章达到真正的统一。前面几章中，我们先后学习了库仑定律、安培力、法拉第电磁感应与交流电路。它们最初看上去像是彼此独立的规律：电荷产生电场，电流产生磁场，变化的磁场又会激发电场。然而麦克斯韦发现，这些规律并不是零散的经验结论，而是同一套场论结构的不同侧面。只要把“位移电流”这一关键概念补入安培定律，电与磁便首尾相接，形成自洽闭环，并自然预言了电磁波的存在。

这正是经典场论最壮丽的一刻：光不再是神秘的“发光物质”，而是电场与磁场彼此激发、共同传播的波动。于是，静电学、稳恒电流、磁场、电磁感应与光学，第一次在一章之中汇合成统一语言。

\section{位移电流——麦克斯韦的关键洞察}

法拉第告诉我们：变化的磁场能够产生旋电场。麦克斯韦进一步追问：如果电磁学真的具有对称的统一结构，那么变化的电场是否也应当产生磁场？这个问题并非纯粹审美，而是由一个具体矛盾逼出来的。

考虑电容器充电过程。导线中有传导电流 $I$ 流向极板，因此按安培环路定律，回路周围应有磁场：
\[
\oint_{\partial \Sigma} \vb{B}\cdot \dd \vb{l}=\mu_0 I_{\text{enc}}.
\]
但若我们对同一条闭合边界 $\partial \Sigma$ 选择两种不同的张面 $\Sigma$：
\begin{itemize}
  \item 一种穿过导线，可截到传导电流 $I$；
  \item 另一种鼓起穿过电容极板间空隙，则没有电子真正跨过间隙。
\end{itemize}
若仍把右端只写成传导电流，便会得到相互矛盾的结果。磁场不可能依赖于我们任意选择的数学曲面，因此原定律一定还缺少一项。

\begin{definition}[位移电流]
位移电流是由电位移矢量 $\vb{D}$ 的时间变化所对应的“等效电流”项，其面通量形式定义为
\[
I_d\coloneqq \dv{}{t}\!\left(\int_{\Sigma} \vb{D}\cdot \dd \vb{A}\right).
\]
若介质为真空，$\vb{D}=\varepsilon_0\vb{E}$，则
\[
I_d=\varepsilon_0\dv{}{t}\!\left(\int_{\Sigma} \vb{E}\cdot \dd \vb{A}\right).
\]
在局域形式中，对应的电流密度称为位移电流密度：
\[
\vb{J}_d=\pdv{\vb{D}}{t}.
\]
\end{definition}

这里“位移电流”并不意味着真的有电荷穿过真空缝隙流动，而是说\emph{变化的电场在磁场方程中起到与电流同样的源作用}。它修补的不是语义，而是守恒律。

\begin{law}[安培--麦克斯韦定律的积分形式]
完整的环路定律应写成
\[
\oint_{\partial \Sigma} \vb{H}\cdot \dd \vb{l}
= I_{\text{free}}+\dv{}{t}\!\left(\int_{\Sigma}\vb{D}\cdot \dd\vb{A}\right).
\]
在真空中，由 $\vb{B}=\mu_0\vb{H}$ 与 $\vb{D}=\varepsilon_0\vb{E}$ 得
\[
\oint_{\partial \Sigma} \vb{B}\cdot \dd \vb{l}
=\mu_0 I_{\text{cond}}+\mu_0\varepsilon_0\dv{}{t}\!\left(\int_{\Sigma}\vb{E}\cdot \dd\vb{A}\right).
\]
\end{law}

为了看清它与电荷守恒的关系，我们回忆连续性方程：
\[
\div \vb{J}+\pdv{\rho}{t}=0.
\]
如果磁场源只写成 $\vb{J}$，那么对方程 $\curl \vb{B}=\mu_0\vb{J}$ 两边取散度，会得到
\[
0=\mu_0\div \vb{J},
\]
这要求 $\div \vb{J}=0$，与一般充放电过程不相容。麦克斯韦把右端补成
\[
\curl \vb{B}=\mu_0\vb{J}+\mu_0\pdv{\vb{D}}{t},
\]
再取散度便有
\[
0=\mu_0\div\vb{J}+\mu_0\pdv{}{t}(\div\vb{D}).
\]
利用高斯定律 $\div\vb{D}=\rho$，立刻恢复连续性方程。这说明位移电流不是可有可无的修辞，而是使场方程与电荷守恒完全兼容的必要项。

\begin{example}[充电电容器中的位移电流]
设平行板电容器极板面积为 $A$，板间距离远小于极板尺度，可近似视为匀强电场。若导线中的传导电流为 $I$，则极板上电荷满足 $\dv{Q}{t}=I$，板间电场为
\[
E=\frac{Q}{\varepsilon_0 A}.
\]
因此位移电流为
\[
I_d=\varepsilon_0\dv{(EA)}{t}=\varepsilon_0 A \dv{E}{t}
=\dv{Q}{t}=I.
\]
这正说明：虽然板间没有传导电流穿过，但位移电流恰好补足了磁场的源，使任意张面得到相同结果。
\end{example}

\begin{figure}[htbp]
\centering
\begin{tikzpicture}
\begin{axis}[
  width=0.75\textwidth, height=0.5\textwidth,
  xlabel={时间 $t/\tau$}, ylabel={位移电流密度 $J_d/J_0$},
  xmin=0, xmax=5, ymin=0, ymax=1.05,
  axis lines=middle,
  grid=major, grid style={gray!30},
  thick,
]
\addplot[blue, domain=0:5, samples=100] {exp(-x)};
\end{axis}
\end{tikzpicture}
\caption{电容器充电过程中位移电流密度随时间指数衰减}
\end{figure}

\begin{remark}
从对称性看，法拉第定律告诉我们“变磁生电”，安培--麦克斯韦定律则告诉我们“变电生磁”。这两个“旋度方程”共同构成电磁波传播的动力学核心。没有位移电流，电磁学只能描述静态与准静态现象；加入位移电流后，场本身便获得脱离物质源而独立传播的能力。
\end{remark}

\begin{center}
\begin{tikzpicture}[scale=1.0, >=Latex]
  % plates
  \draw[fill=blue!8] (-3,1.8) rectangle (-2,-1.8);
  \draw[fill=red!8] (2,1.8) rectangle (3,-1.8);
  \node at (-2.5,2.2) {负极板};
  \node at (2.5,2.2) {正极板};

  % wires
  \draw[thick] (-5,0.8) -- (-3,0.8);
  \draw[thick] (3,0.8) -- (5,0.8);
  \draw[->, thick] (-4.4,0.8) -- (-3.3,0.8);
  \node at (-4.2,1.2) {$I$};

  % E field
  \foreach \y in {-1.2,-0.4,0.4,1.2} {
    \draw[->, red!70!black, thick] (1.5,\y) -- (-1.5,\y);
  }
  \node[red!70!black] at (0,1.7) {$\vb{E}(t)$ 增强};

  % displacement current
  \draw[->, dashed, blue!80!black, very thick] (0,-1.4) -- (0,1.4);
  \node[blue!80!black, right] at (0.1,0) {$\pdv{\vb{D}}{t}$};

  % magnetic loops
  \draw[orange!80!black, thick, ->] (0,0) ellipse (1.3 and 2.0);
  \node[orange!80!black] at (1.5,-1.8) {$\vb{B}$ 环绕};
\end{tikzpicture}
\end{center}

\section{麦克斯韦方程组——统一的四个方程}

位移电流一旦加入，经典电磁学便凝聚为四个基本方程。它们一方面说明电荷与电流如何产生场，另一方面说明变化的场如何彼此激发。这四个方程既可写成整体的积分形式，也可写成逐点成立的微分形式。

\begin{law}[麦克斯韦方程组的积分形式]
在宏观介质表述中，麦克斯韦方程组为
\begin{align*}
\oint_{\partial V} \vb{D}\cdot \dd \vb{A} &= Q_{\text{free, enc}}, \\
\oint_{\partial V} \vb{B}\cdot \dd \vb{A} &= 0, \\
\oint_{\partial \Sigma} \vb{E}\cdot \dd \vb{l} &= -\dv{}{t}\!\left(\int_{\Sigma}\vb{B}\cdot \dd \vb{A}\right), \\
\oint_{\partial \Sigma} \vb{H}\cdot \dd \vb{l} &= I_{\text{free, enc}}+\dv{}{t}\!\left(\int_{\Sigma}\vb{D}\cdot \dd \vb{A}\right).
\end{align*}
\end{law}

它们依次对应：
\begin{enumerate}
  \item 电场的源是自由电荷；
  \item 不存在磁单极子，磁力线总是闭合；
  \item 变化的磁通量激发环形电场；
  \item 自由电流与变化的电通量共同激发环形磁场。
\end{enumerate}

在真空中使用 $\vb{D}=\varepsilon_0\vb{E}$、$\vb{B}=\mu_0\vb{H}$，方程可化为更熟悉的形式：
\begin{align*}
\oint_{\partial V} \vb{E}\cdot \dd \vb{A} &= \frac{Q_{\text{enc}}}{\varepsilon_0}, \\
\oint_{\partial V} \vb{B}\cdot \dd \vb{A} &= 0, \\
\oint_{\partial \Sigma} \vb{E}\cdot \dd \vb{l} &= -\dv{}{t}\!\left(\int_{\Sigma}\vb{B}\cdot \dd \vb{A}\right), \\
\oint_{\partial \Sigma} \vb{B}\cdot \dd \vb{l} &= \mu_0 I_{\text{enc}}+\mu_0\varepsilon_0\dv{}{t}\!\left(\int_{\Sigma}\vb{E}\cdot \dd \vb{A}\right).
\end{align*}

利用高斯定理与斯托克斯定理，积分形式可转化为局域微分形式。

\begin{theorem}[麦克斯韦方程组的微分形式]
在真空中，麦克斯韦方程组可写为
\begin{align*}
\div \vb{E} &= \frac{\rho}{\varepsilon_0}, \\
\div \vb{B} &= 0, \\
\curl \vb{E} &= -\pdv{\vb{B}}{t}, \\
\curl \vb{B} &= \mu_0\vb{J}+\mu_0\varepsilon_0\pdv{\vb{E}}{t}.
\end{align*}
在一般介质中，则写成
\begin{align*}
\div \vb{D} &= \rho_{\text{free}}, \\
\div \vb{B} &= 0, \\
\curl \vb{E} &= -\pdv{\vb{B}}{t}, \\
\curl \vb{H} &= \vb{J}_{\text{free}}+\pdv{\vb{D}}{t}.
\end{align*}
\end{theorem}

\begin{proof}
以法拉第定律为例，积分形式为
\[
\oint_{\partial \Sigma}\vb{E}\cdot \dd\vb{l}=-\dv{}{t}\!\left(\int_\Sigma \vb{B}\cdot \dd\vb{A}\right).
\]
由斯托克斯定理，左端可写为
\[
\int_\Sigma (\curl \vb{E})\cdot \dd\vb{A}.
\]
又因该关系对任意小曲面 $\Sigma$ 成立，所以被积函数必须逐点相等，故
\[
\curl \vb{E}=-\pdv{\vb{B}}{t}.
\]
其余三个方程的推导完全类似：对体积分方程使用高斯定理，对环路积分方程使用斯托克斯定理即可。
\end{proof}

\begin{definition}[真空中的电磁场]
若某区域内没有自由电荷与传导电流，即
\[
\rho=0,\qquad \vb{J}=0,
\]
则该区域满足真空麦克斯韦方程组
\begin{align*}
\div \vb{E} &= 0, & \div \vb{B} &= 0,\\
\curl \vb{E} &= -\pdv{\vb{B}}{t}, & \curl \vb{B} &= \mu_0\varepsilon_0\pdv{\vb{E}}{t}.
\end{align*}
这组齐次方程正是电磁波的起点。
\end{definition}

\begin{example}[平行板电容器边缘处的磁场]
设一半径为 $r$ 的圆形积分回路位于平行板电容器内部且与极板同轴，若板间电场近似均匀，则电通量为
\[
\Phi_E=E\pi r^2.
\]
安培--麦克斯韦定律给出
\[
\oint \vb{B}\cdot \dd\vb{l}=B(2\pi r)=\mu_0\varepsilon_0 \dv{\Phi_E}{t}
=\mu_0\varepsilon_0 \pi r^2 \dv{E}{t}.
\]
因此
\[
B=\frac{\mu_0\varepsilon_0 r}{2}\dv{E}{t}.
\]
它表明即使在没有传导电流穿过的区域，只要电场随时间变化，磁场依然存在。
\end{example}

\begin{remark}
从结构上看，前两个散度方程描述“源”的分布，后两个旋度方程描述“变化”如何驱动传播。静电学和稳恒磁场只是这套理论在时间不变极限下的两个投影；完整的电磁学则是场在时空中耦合演化的理论。
\end{remark}

\section{电磁波的推导——场如何自我传播}

一旦进入真空区域，电荷和电流都消失，场方程反而变得更有力量。因为此时电场与磁场不再依赖外部源维持，而是通过旋度方程彼此驱动，自行向前传播。

在真空中，麦克斯韦方程组为
\begin{align}
\div \vb{E} &= 0, \label{eq:maxwell-vacuum-divE}\\
\div \vb{B} &= 0, \label{eq:maxwell-vacuum-divB}\\
\curl \vb{E} &= -\pdv{\vb{B}}{t}, \label{eq:maxwell-vacuum-curlE}\\
\curl \vb{B} &= \mu_0\varepsilon_0\pdv{\vb{E}}{t}. \label{eq:maxwell-vacuum-curlB}
\end{align}
对 \cref{eq:maxwell-vacuum-curlE} 两边取旋度：
\[
\curl(\curl \vb{E})=-\pdv{}{t}(\curl \vb{B}).
\]
代入 \cref{eq:maxwell-vacuum-curlB}，得
\[
\curl(\curl \vb{E})=-\mu_0\varepsilon_0\pdv[2]{\vb{E}}{t}.
\]
再用向量恒等式
\[
\curl(\curl \vb{E})=\grad(\div\vb{E})-\nabla^2 \vb{E},
\]
并结合 $\div\vb{E}=0$，便得到
\[
\nabla^2 \vb{E}=\mu_0\varepsilon_0\pdv[2]{\vb{E}}{t}.
\]
这正是三维波动方程。对磁场重复同样步骤，可得
\[
\nabla^2 \vb{B}=\mu_0\varepsilon_0\pdv[2]{\vb{B}}{t}.
\]

\begin{theorem}[真空中的电磁波动方程]
在无源真空区域内，电场与磁场分别满足
\[
\nabla^2 \vb{E}=\mu_0\varepsilon_0\pdv[2]{\vb{E}}{t},
\qquad
\nabla^2 \vb{B}=\mu_0\varepsilon_0\pdv[2]{\vb{B}}{t}.
\]
因此电磁扰动以波速
\[
 v=\frac{1}{\sqrt{\mu_0\varepsilon_0}}
\]
传播。
\end{theorem}

若代入实验测得的真空磁导率与介电常数，便得到
\[
\frac{1}{\sqrt{\mu_0\varepsilon_0}}\approx 3.00\times 10^8\ \text{m/s},
\]
这恰好等于光速 $c$。因此麦克斯韦得出震撼性的结论：
\begin{center}
\Large 光就是电磁波。
\end{center}

这是理论物理史上的经典时刻：不是先看到一种现象再去拟合方程，而是先从方程中推出一个此前独立研究的自然对象，并发现它与已知常数严丝合缝地吻合。

\begin{example}[一维平面电磁波满足波动方程]
设电场仅沿 $y$ 方向，且只依赖于 $x,t$：
\[
\vb{E}(x,t)=E_y(x,t)\,\uvec{y}.
\]
则波动方程化为
\[
\pdv[2]{E_y}{x}=\mu_0\varepsilon_0\pdv[2]{E_y}{t}.
\]
其一组解为
\[
E_y(x,t)=E_0\cos(kx-\omega t+\phi),
\]
只要满足色散关系
\[
\omega=ck,\qquad c=\frac{1}{\sqrt{\mu_0\varepsilon_0}}.
\]
磁场也可取同频同相的平面波形式，从而组成完整的电磁波解。
\end{example}

\begin{law}[平面电磁波的场关系]
对沿 $+x$ 方向传播的简谐平面波，可取
\[
\vb{E}=E_0\cos(kx-\omega t)\,\uvec{y},
\qquad
\vb{B}=B_0\cos(kx-\omega t)\,\uvec{z}.
\]
代入法拉第定律或安培--麦克斯韦定律可得
\[
B_0=\frac{E_0}{c},
\qquad
\omega=ck.
\]
\end{law}

\begin{figure}[htbp]
\centering
\begin{tikzpicture}
\begin{axis}[
  width=0.75\textwidth, height=0.5\textwidth,
  xlabel={位置 $x/\lambda$}, ylabel={归一化场强},
  xmin=0, xmax=2, ymin=-1.2, ymax=1.2,
  axis lines=middle,
  grid=major, grid style={gray!30},
  thick,
  legend style={at={(0.98,0.02)}, anchor=south east, draw=none, fill=none}
]
\addplot[blue, domain=0:2, samples=200] {sin(deg(2*pi*x))};
\addlegendentry{$E/E_0$}
\addplot[red, dashed, domain=0:2, samples=200] {sin(deg(2*pi*x))};
\addlegendentry{$cB/E_0$}
\end{axis}
\end{tikzpicture}
\caption{沿传播方向的电场与磁场同相作正弦振荡}
\end{figure}

\section{电磁波的性质——横波、正交与能量}

麦克斯韦方程组不仅告诉我们电磁波会传播，还严格规定了它如何传播。与绳波、声波等机械波不同，电磁波不需要介质振动，其“振动”对象就是场本身。

\subsection*{1. 横波性质}

设平面波沿波矢 $\vb{k}$ 方向传播，场可写成
\[
\vb{E}(\vb{r},t)=\vb{E}_0\cos(\vb{k}\cdot\vb{r}-\omega t),
\qquad
\vb{B}(\vb{r},t)=\vb{B}_0\cos(\vb{k}\cdot\vb{r}-\omega t).
\]
由 $\div\vb{E}=0$ 与 $\div\vb{B}=0$ 可得
\[
\vb{k}\cdot\vb{E}_0=0,
\qquad
\vb{k}\cdot\vb{B}_0=0.
\]
因此电场与磁场都垂直于传播方向，电磁波是\textbf{横波}。

\subsection*{2. $\vb{E}\perp \vb{B}$ 且同相}

由旋度方程进一步可知
\[
\vb{k}\times \vb{E}_0=\omega \vb{B}_0,
\qquad
\vb{k}\times \vb{B}_0=-\frac{\omega}{c^2}\vb{E}_0.
\]
因此 $\vb{E}_0\perp \vb{B}_0$，且 $\vb{E},\vb{B},\vb{k}$ 三者两两正交，构成右手系。对真空平面波而言，电场与磁场在相位上同步达到极大值与极小值，不存在四分之一周期的相位差。

\subsection*{3. 振幅关系与能量均分}

真空中平面波满足
\[
E=cB.
\]
电场能量密度与磁场能量密度分别为
\[
u_E=\frac{1}{2}\varepsilon_0 E^2,
\qquad
u_B=\frac{1}{2\mu_0}B^2.
\]
利用 $E=cB$ 以及 $c^2=1/(\mu_0\varepsilon_0)$，立刻得到
\[
u_E=u_B.
\]
所以在真空平面电磁波中，电场与磁场各承担一半能量。总能量密度为
\[
u=u_E+u_B=\frac{1}{2}\varepsilon_0 E^2+\frac{B^2}{2\mu_0}.
\]

\begin{definition}[电磁波的能量密度]
电磁场在真空中的瞬时能量密度定义为
\[
 u\coloneqq \frac{1}{2}\varepsilon_0 E^2+\frac{B^2}{2\mu_0}.
\]
对于简谐平面波，由于 $E=cB$，可写成
\[
 u=\varepsilon_0 E^2=\frac{B^2}{\mu_0}.
\]
若对一个周期取时间平均，则
\[
 \avg{u}=\frac{1}{2}\varepsilon_0 E_0^2=\frac{B_0^2}{2\mu_0}.
\]
\end{definition}

\begin{definition}[偏振]
若一列沿固定方向传播的电磁波，其电场矢量在垂直于传播方向的平面内具有确定的振动规律，则称该波具有偏振。若 $\vb{E}$ 始终沿固定方向振动，称为线偏振；若其端点描绘圆或椭圆，则分别称为圆偏振或椭圆偏振。
\end{definition}

偏振现象是电磁波横波性的直接证据之一。自然光可视作许多偏振态的混合，而偏振片则会筛选出某一方向的电场分量。若电磁波是纵波，就不会存在这样的“方向选择”现象。

\begin{example}[太阳光辐照下的场强估算]
设太阳光到达地面的平均辐照度约为 $I=1.0\times 10^3\ \text{W/m}^2$。对平面电磁波，平均坡印廷矢量大小满足
\[
I=\avg{S}=\frac{1}{2}\varepsilon_0 c E_0^2.
\]
因此
\[
E_0=\sqrt{\frac{2I}{\varepsilon_0 c}}\approx 8.7\times10^2\ \text{V/m}.
\]
相应磁场振幅为
\[
B_0=\frac{E_0}{c}\approx 2.9\times10^{-6}\ \text{T}.
\]
这说明日常光照对应的磁场其实非常弱，但其高速传播使能量输运极其显著。
\end{example}

\begin{center}
\begin{tikzpicture}[scale=1.0, >=Latex]
  \draw[->] (0,0) -- (7.2,0) node[right] {$x$（传播方向）};
  \draw[->, red!75!black] (0,0) -- (0,2.5) node[above] {$E_y$};
  \draw[->, blue!75!black] (0,0) -- (0,-2.5) node[below] {$B_z$};

  \draw[red!75!black, thick, samples=120, domain=0:6.3, smooth]
    plot (\x,{1.4*sin(\x r)});
  \draw[blue!75!black, thick, dashed, samples=120, domain=0:6.3, smooth]
    plot (\x,{-1.2*sin(\x r)});

  \draw[->, very thick, orange!80!black] (3.8,0.3) -- (5.3,0.3);
  \node[orange!80!black] at (4.55,0.8) {$\vb{S}$};
  \node[red!75!black] at (6.1,1.5) {$\vb{E}$};
  \node[blue!75!black] at (6.0,-1.4) {$\vb{B}$};
\end{tikzpicture}
\end{center}

\begin{remark}
声波依赖空气的压缩与疏松，所以不能在真空中传播；电磁波则不同，电场与磁场互为“弹性”和“惯性”，场本身就构成传播机制。这是经典物理从“粒子受力”走向“场独立存在”的一次哲学跃迁。
\end{remark}

\section{坡印廷矢量——电磁能如何流动}

既然电磁波能够携带能量，那么能量流向何处、以何种速率穿过单位面积，就需要新的物理量来描述。这个量就是坡印廷矢量。

\begin{definition}[坡印廷矢量]
真空中电磁场的能流密度矢量定义为
\[
\vb{S}\coloneqq \frac{1}{\mu_0}\vb{E}\times\vb{B}.
\]
其方向给出电磁能传播的方向，其大小表示单位时间穿过单位面积的能量，即功率面密度。
\end{definition}

对平面波而言，由于 $\vb{E}\perp\vb{B}$，有
\[
S=\frac{EB}{\mu_0}=\varepsilon_0 c E^2=\frac{c}{\mu_0}B^2.
\]
若电场为简谐振荡 $E=E_0\cos(kx-\omega t)$，则瞬时能流密度也随时间振荡：
\[
S(t)=\varepsilon_0 c E_0^2 \cos^2(kx-\omega t).
\]
一个周期内的平均值为
\[
\avg{S}=\frac{1}{2}\varepsilon_0 c E_0^2.
\]
这就是光强、无线电发射强度等量的经典表达。

\begin{theorem}[坡印廷定理的真空形式]
设真空中存在电荷密度 $\rho$ 和电流密度 $\vb{J}$，则电磁能满足局域守恒关系
\[
\pdv{u}{t}+\div \vb{S}+\vb{J}\cdot\vb{E}=0,
\]
其中
\[
 u=\frac{1}{2}\varepsilon_0 E^2+\frac{B^2}{2\mu_0},
 \qquad
 \vb{S}=\frac{1}{\mu_0}\vb{E}\times\vb{B}.
\]
\end{theorem}

\begin{proof}
由安培--麦克斯韦定律
\[
\curl\vb{B}=\mu_0\vb{J}+\mu_0\varepsilon_0\pdv{\vb{E}}{t}
\]
与法拉第定律
\[
\curl\vb{E}=-\pdv{\vb{B}}{t}
\]
出发，分别与 $\vb{E}$、$\vb{B}/\mu_0$ 做点乘并相减，利用向量恒等式
\[
\div(\vb{E}\times\vb{B})=\vb{B}\cdot(\curl\vb{E})-\vb{E}\cdot(\curl\vb{B}),
\]
整理即可得到
\[
\vb{J}\cdot\vb{E}+\pdv{}{t}\left(\frac{1}{2}\varepsilon_0 E^2+\frac{B^2}{2\mu_0}\right)+\div\left(\frac{1}{\mu_0}\vb{E}\times\vb{B}\right)=0.
\]
这就是所求结果。
\end{proof}

把上式对任意体积 $V$ 积分，并使用高斯定理，可得
\[
\dv{}{t}\!\left(\int_V u\,\dd V\right)+\oint_{\partial V}\vb{S}\cdot\dd\vb{A}+\int_V \vb{J}\cdot\vb{E}\,\dd V=0.
\]
它表明：边界面上的坡印廷通量给出了体积内部电磁能的净流出率，而 $\vb{J}\cdot\vb{E}$ 则给出了场能转化为物质能的速率。

坡印廷定理的物理意义极其清晰：
\begin{itemize}
  \item $\pdv{u}{t}$ 表示单位体积内电磁能的变化率；
  \item $\div\vb{S}$ 表示能量从该体积流出或流入；
  \item $\vb{J}\cdot\vb{E}$ 表示电磁场对电荷做功、把场能转为粒子能量的速率。
\end{itemize}
因此它是电磁场的能量守恒律。

作为数量级估算，若某真空平面电磁波的电场振幅为 $E_0=120\ \text{V/m}$，则平均能流密度为
\[
\avg{S}=\frac{1}{2}\varepsilon_0 c E_0^2
\approx \frac{1}{2}\times 8.85\times10^{-12}\times 3.00\times10^8\times (120)^2
\approx 19\ \text{W/m}^2.
\]
这表示每秒约有 $19\ \text{J}$ 的电磁能穿过垂直于传播方向的每平方米面积。

\begin{remark}
在稳恒直流电路中，我们常说电源通过导线把能量送到电阻上。更精确地说，真正携带能量的不是电子“拖着能量跑”，而是导线外部空间中的电磁场。电场大致沿导线方向，磁场环绕导线，因此
\[
\vb{S}=\frac{1}{\mu_0}\vb{E}\times\vb{B}
\]
指向导线内部或电阻内部。于是，能量是从周围空间持续流入电阻并转化为热的。这是坡印廷矢量给出的深刻图景。
\end{remark}

\begin{remark}
“光压”“太阳帆”“无线能量传输”等现象，都可以从坡印廷矢量出发理解。它把“场携带能量与动量”这一观点变成可计算、可测量的物理量，也为后续进入相对论与场论铺平道路。
\end{remark}

\section{电磁谱——从无线电波到 $\gamma$ 射线}

既然光只是电磁波的一小段，那么不同频率、不同波长的所有电磁辐射，本质上都由同一套麦克斯韦方程支配。改变的不是“种类”，而只是频率尺度与相互作用方式。于是，人类把分散的无线通信、红外加热、可见光成像、紫外灭菌、X 射线透视、$\gamma$ 射线核辐射统统纳入一个连续谱中。

\begin{definition}[电磁谱]
按照频率 $f$ 或波长 $\lambda$ 排列的全部电磁波范围称为电磁谱，它们满足
\[
 c=f\lambda.
\]
频率越高，波长越短；在真空中它们的传播速度都等于光速 $c$。
\end{definition}

\begin{center}
\renewcommand{\arraystretch}{1.25}
\begin{tabular}{p{3cm}p{3.7cm}p{4.8cm}p{3.2cm}}
\toprule
波段 & 典型波长范围 & 典型应用 & 主要特征 \\
\midrule
无线电波 & $>10^{-1}\ \text{m}$ & 广播、通信、雷达 & 波长长，绕射能力强 \\
微波 & $10^{-3}\sim10^{-1}\ \text{m}$ & 微波炉、卫星通信、Wi-Fi & 易与分子转动耦合 \\
红外线 & $7\times10^{-7}\sim10^{-3}\ \text{m}$ & 热成像、遥控 & 常与热辐射关联 \\
可见光 & $4\times10^{-7}\sim7\times10^{-7}\ \text{m}$ & 视觉、成像、光纤 & 人眼可感知部分 \\
紫外线 & $10^{-8}\sim4\times10^{-7}\ \text{m}$ & 灭菌、荧光激发 & 化学作用强 \\
X 射线 & $10^{-11}\sim10^{-8}\ \text{m}$ & 医学透视、晶体分析 & 穿透性强 \\
$\gamma$ 射线 & $<10^{-11}\ \text{m}$ & 核物理、放疗、天体物理 & 频率最高、能量最大 \\
\bottomrule
\end{tabular}
\end{center}

\begin{example}[由频率求波长]
某 FM 广播频率为 $f=100\ \text{MHz}$，则其在真空中的波长为
\[
\lambda=\frac{c}{f}=\frac{3.00\times 10^8}{1.00\times10^8}=3.0\ \text{m}.
\]
这说明广播电台使用的是米量级波长的无线电波。
\end{example}

若可见光典型波长取 $500\ \text{nm}$，X 射线典型波长取 $0.1\ \text{nm}$，则二者波长之比为
\[
\frac{500}{0.1}=5000.
\]
也就是说，X 射线的波长比可见光短约四个数量级，因此更容易分辨原子尺度结构，也具有更强穿透能力。

\begin{remark}
麦克斯韦方程组本身并不偏爱任何波段。无线电通信与宇宙射线看似相距遥远，但它们只是同一种场在不同频率下的表现。统一的力量，就在于把“多样现象”还原为“少数基本方程”。从经典观点看，不同波段的差异主要体现在与物质的耦合尺度不同；从现代观点看，还可以进一步用光子的能量 $E=hf$ 来理解高频辐射为何更容易电离原子。
\end{remark}

\section*{本章小结}

本章是电磁学的顶点，也是经典物理统一思想的典范：
\begin{itemize}
  \item 麦克斯韦以位移电流 $\pdv{\vb{D}}{t}$ 修正安培定律，使其与电荷守恒完全一致；
  \item 四条麦克斯韦方程把电场、磁场、电荷和电流纳入统一框架；
  \item 在真空中，这四个方程自动推出波动方程，预言波速 $c=1/\sqrt{\mu_0\varepsilon_0}$；
  \item 电磁波是横波，满足 $\vb{E}\perp \vb{B}\perp \vb{k}$，并携带能量与动量；
  \item 坡印廷矢量 $\vb{S}=\vb{E}\times\vb{B}/\mu_0$ 描述了电磁能流；
  \item 从无线电波到 $\gamma$ 射线，整个电磁谱都服从同一组方程。
\end{itemize}

如果说牛顿把天上与地下的运动统一起来，那么麦克斯韦则把电、磁与光统一起来。下一步，爱因斯坦将进一步追问：既然光速由真空常数决定，而且对所有惯性系都应相同，那么时间与空间本身又必须如何改变？

